\section{Introduction}
\label{sec:introduction}

Large language models (LLMs) are now strong zero-shot rerankers in
information retrieval (IR) pipelines. The dominant paradigms differ in how
they elicit relevance judgements from the model: pointwise scoring evaluates
one query-document pair at a time~\citep{nogueira2020monot5}, pairwise
prompting compares two documents per
call~\citep{qin2024prp}, listwise prompting feeds an entire candidate list and
asks for a permutation~\citep{sun2023rankgpt,ma2023zsl}, and
\emph{setwise} prompting~\citep{zhuang2024setwise} feeds a small candidate
window of $c{+}1$ passages and asks the model to pick the most relevant one,
embedding that comparator inside a classical sort (heapsort, bubblesort) to
produce a top-$k$ ranking. Setwise reranking has become a standard
zero-shot baseline because it sits on the effectiveness--efficiency Pareto
frontier between pointwise and listwise, while keeping prompts short enough to
fit comfortably in any modern LLM's context.

Yet the setwise comparator is informationally thin. Each LLM call reads
$c{+}1$ passages, internally evaluates all of them, and emits a single label.
The algorithm uses only that one explicit extreme --- the identity of the best
document --- and does not explicitly use the rest of the model's implicit
comparison evidence. Two recent directions try to use each setwise call more
effectively. Joint elicitation asks for both the best and worst document in
one prompt (\textsc{DualEnd}), recovering two explicit extremes per call.
Tournament prompting~\citep{chen2025tourrank} keeps the best-document
comparator but embeds it in a bracket, so advancement, elimination, and
accumulated points turn local winners into broader ranking evidence. Both
directions improve the information-per-call envelope, but both remain
anchored to bounded local groups rather than repeatedly prompting over the
whole rerank pool.

\paragraph{The mismatch.} Modern instruction-tuned LLMs comfortably handle
contexts of 32K--256K tokens. The Qwen3 family used in this paper, for
instance, ships with 32K native context for the 4B/8B/14B variants and 256K
for the Qwen3.5 9B model. At a passage length of $512$ tokens, a $50$-document
rerank pool plus query, scaffolding, and decoding budget consumes well under
$30\mathrm{K}$ tokens --- it fits in one prompt. The setwise sliding window of
size $c{+}1{=}4$ is therefore an artefact of an earlier compute regime: it
forces hundreds of sequential LLM calls per query whose only role is to
compose enough small local decisions into a top-$k$ ranking, even though the
entire rerank pool would fit in a single context window. In our completed
baseline runs on TREC DL19/20, \textsc{DualEnd-Cocktail} averages $546$
comparisons and roughly $8.9{\times}$ the wall-clock of a TopDown-Heap
baseline, for a mean NDCG@10 gain of only $+0.0065$ over the cheapest setwise
sort. The Pareto frontier between TopDown-Bubble (300 comparisons, 110.9~s,
NDCG@10 $0.6897$) and DualEnd-Cocktail (546 comparisons, 212.6~s, NDCG@10
$0.6962$) is \emph{empty}: none of our completed setwise / DualEnd baselines
delivers DualEnd-class quality at TopDown-class cost.

\paragraph{This paper.} We target this gap by abandoning the small-window
assumption altogether. We introduce \textbf{MaxContext}, a family of zero-shot
LLM rerankers that prompts the model over the \emph{entire live rerank pool}
on every call. Given a starting pool of $N$ documents (we evaluate
$N \in \{10, 20, 30, 40, 50\}$, all bounded by $50$), MaxContext repeatedly
sends the live pool to the LLM, removes the document(s) the model points to,
and continues until the pool is exhausted. Three variants probe different
hypotheses about \emph{which} signal the long-context model should produce on
each call:
\begin{itemize}
\item \textsc{MaxContext-DualEnd} asks for both the best and worst document
  in the live pool; the pool shrinks by two documents per call. Total cost:
  $\lfloor N/2 \rfloor$ LLM calls per query (e.g.\ $25$ at $N{=}50$).
\item \textsc{MaxContext-TopDown} asks for only the best document; the pool
  shrinks by one. A two-document deterministic BM25 endgame resolves the last
  pair. Total cost: $N{-}2$ calls plus a constant-time tiebreaker.
\item \textsc{MaxContext-BottomUp} asks for only the worst document; the pool
  shrinks by one with the same BM25 endgame. Total cost matches TopDown but
  on the opposite end of the ranking.
\end{itemize}
Each call places one or two documents at globally observed rank positions, so
the algorithm produces a permutation of the input pool with no early
termination, no auto-batching, and no auto-truncation. We replace the legacy
letter labels (\texttt{A}--\texttt{W}) with numeric labels $1$--$N$ to support
$N{\,>\,}23$, harden the parser against silent fallbacks at $N{=}50$, and
introduce strict invariants ($\texttt{pool\_size} = \texttt{hits} =
\texttt{ranker.k}$, abort-on-bad-parse, no truncation) so that whole-pool
prompts cannot be silently clipped or misinterpreted.

\paragraph{Why three variants.} MaxContext is a method family rather than a
single method: it implements a controlled comparison of comparator output
types under a fixed whole-pool prompt structure. Three predeclared hypotheses
are testable simultaneously:
\begin{enumerate}
\item \textbf{Joint-elicitation hypothesis.} Whole-pool quality comes from
  asking for two extremes per call. Predicts
  \textsc{DualEnd}~$\succ$~\textsc{TopDown}~$\approx$~\textsc{BottomUp}.
\item \textbf{Whole-pool-selection hypothesis.} Whole-pool quality comes from
  letting the model see all candidates simultaneously, regardless of how many
  outputs it emits. Predicts
  \textsc{DualEnd}~$\approx$~\textsc{TopDown}~$\approx$~\textsc{BottomUp}.
\item \textbf{Single-extreme hypothesis.} Joint elicitation does not scale to
  long context and degrades \textsc{DualEnd} relative to selection-only
  variants. Predicts one or both single-extreme variants outperform
  \textsc{DualEnd}.
\end{enumerate}
Reporting all three variants under a common harness lets the empirical
outcome itself decide which mechanism is responsible for any frontier
movement, rather than committing in advance to one explanation.

\paragraph{Complementary operating point: Top-$m$ + DualEnd cascade.}
MaxContext's hard ceiling at $\text{pool size}\le 50$ is set by the smallest
context window in our model lineup. To probe the same Pareto-frontier region
\emph{without} a hard pool-size cap, we additionally evaluate a Top-$m$
Setwise + DualEnd finalist cascade. The cascade asks the LLM to pick the
top-$m$ documents from each setwise group, accumulates votes across $r$
stratified groupings, and then refines the surviving finalist pool with
MaxContext-DualEnd. The cascade can in principle scale to larger initial
pools at the cost of weaker recall guarantees, and so serves as a different
operating point on the same comparisons/wall-clock frontier rather than as a
competing main method.

\paragraph{Reporting discipline.} A long-context whole-pool prompt is
\emph{not} a free lunch. It uses substantially more prompt tokens per call
than a $4$-passage setwise prompt; the total number of dependent LLM calls
can drop by an order of magnitude (e.g., for MaxContext-DualEnd relative to
small-window DualEnd-Cocktail), but per-call token cost rises in tandem. We
therefore commit, throughout the paper, to reporting all five cost axes
recommended by the recent reranker efficiency
literature~\citep{peng2025flopsreranking}: total dependent LLM calls
(``comparisons''), prompt tokens, completion tokens, total tokens, and
end-to-end wall-clock latency. Our efficiency claim is framed strictly as
\emph{movement on the comparisons-axis and wall-clock-axis}, not as raw token
efficiency. A reader interested in pure token cost will find \textsc{TD-Heap}
remains the cheapest option in our experiments; MaxContext is designed for
readers interested in sequential-call count or end-to-end latency at
competitive NDCG@10 quality.

\paragraph{Contributions.}
\begin{enumerate}
\item \textbf{Method.} We introduce MaxContext, a zero-shot setwise LLM
  reranking family that prompts the model over the entire live rerank pool on
  every call. We give three variants spanning a controlled comparator-output
  ablation (joint best+worst, best-only, worst-only) and present the
  algorithmic, prompt-engineering, and parser-hardening details required to
  make whole-pool prompts robust at $N{=}50$.
\item \textbf{Pareto-frontier characterization.} We empirically characterize
  the previously empty region between TopDown-Bubble and DualEnd-Cocktail on
  the comparisons-axis and wall-clock-axis frontiers, with matched-\texttt{hits}
  baselines at three pool sizes and full transparent cost reporting.
\item \textbf{Mechanism diagnosis.} The three-variant ablation, combined with
  pool-size sweeps and an order-robustness pilot, lets us diagnose
  \emph{whether} any quality difference comes from joint elicitation,
  whole-pool selection itself, or single-extreme selection at long context
  --- a question that no existing setwise study can answer because none
  reaches large enough pools to expose the difference.
\item \textbf{Complementary operating point.} We additionally evaluate a
  Top-$m$ Setwise + DualEnd cascade as a second operating point on the same
  frontier, providing a method that can in principle scale beyond the $50$-
  document pool ceiling.
\item \textbf{Order-robustness and position-bias-at-scale diagnostics.} We
  predeclare a controlled-ordering pilot at $N{=}50$ and report
  selected-position frequencies under a numeric label scheme $1$--$N$,
  separately from the $w{=}4$ letter-scheme position-bias evidence in prior
  work, so that any extension or contradiction of the small-window
  position-bias literature is reported as an empirical finding rather than
  presupposed.
\end{enumerate}

The paper is organized as follows. Section~\ref{sec:related_work} situates
MaxContext among setwise extensions, tournament-style rerankers, and the
long-context-attention literature. Section~\ref{sec:methodology} formalises
MaxContext and the Top-$m$+DualEnd cascade. Section~\ref{sec:experimental_setup}
describes the staged evaluation matrix; Sections~\ref{sec:results}
and~\ref{sec:analysis} report the matched-hits Pareto frontier, the
mechanism-ablation, and the order-robustness diagnostics.
Section~\ref{sec:conclusion} concludes.
